\documentclass[12pt,a4paper]{article}

\setlength{\topmargin}{0.0in}
\setlength{\oddsidemargin}{0.33in}
\setlength{\textheight}{9.0in}
\setlength{\textwidth}{6.0in}
\renewcommand{\baselinestretch}{1.25}



\usepackage{hyperref}
% \usepackage[utf8]{inputenc}
% \usepackage[T1]{fontenc}
\usepackage{float}
\usepackage{enumitem}
\usepackage{amsmath}
\usepackage{amssymb}
\usepackage{tikz}
\usepackage{amsfonts}
\usepackage{graphicx}
\usepackage{subcaption}
\usepackage{caption}
\newcommand{\dv}[2]{\frac{\mathrm{d} #1}{\mathrm{d} #2}}
\usepackage{caption}
\captionsetup[figure]{font=small} 
% \captionsetup[table]{font=small}
\usepackage[capitalize,nameinlink]{cleveref}
\usepackage{mathtools}
\newcommand{\defeq}{\vcentcolon=}
\newcommand{\eqdef}{=\vcentcolon}

% Define environments
\newtheorem{theorem}{Theorem}
\newtheorem{definition}{Definition}
\newtheorem{lemma}{Lemma}
\newtheorem{corollary}{Corollary}



\title{Logs from May $19$ to the $25^{th}$}
\author{}
\date{}
% \author{Maggie McCarthy}
% \date{May 2026}

\begin{document}

\maketitle

\section{Meeting May 19}

\subsection{Summary of last week's work:}




% \begin{enumerate}
%     \item The repo is up and runnning with meeting logs as well as with daily logs of my progress.
%     \item Last week I focused alot on some functional analysis review. This included
%     \begin{itemize}
%         \item Operator review (linear, bounded, closed, compact, inverse, etc).
%         \item Introduction to the spectrum of a linear operator. This included the distinction between the point, continuous, and residual spectra. 
%         \item I learned about Banach and Hilbert adjoints from Umberto's first lecture. 
%         \item Compared self-adjoint, unitary, and normal operators. 
%         \item Learned about the Fredholm alternative for compact linear operators from Umberto's first lecture.
%     \end{itemize}
%     \item Went through the first chapter of Nick's text. This included learning about
%     \begin{itemize}
%         \item The general motivations for computing eigenvalues.
%         \item The motivations for computing pseudospectra of nonnormal operators.
%         \item Eigenvalues of matrices and linear operators.
%         \item Pseudospectra of matrices.
%         \item Pseudospectra of linear operators. 
%     \end{itemize}
%     \item To further consolidate this material, I went through Umberto's first two lectures. Here I learned about 
%     \begin{itemize}
%         \item Hilbert vs. Banach vs. formal adjoints.
%         \item Spectra and pseudospectra of linear operators.
%         \item Spectral radii of bounded operators. 
%         \item Some properties of spectra for different operator-types. For example, the spectrum of a compact operator consists of isolated points.
%         \item Separable Hilbert spaces and Hilbert bases.
%         \item Spectral theorem for operators that are both compact and self-adjoint or normal.
%     \end{itemize}
%     \item Finally, I went through the preliminary section of Matt's text to fill in any gaps. Here, I learned more about
%     \begin{itemize}
%         \item Separable Hilbert spaces and orthogonal projections.
%         \item Uniform boundedness theorem (review).
%         \item His definition of the adjoint. 
%         \item Boundedly invertible operators. 
%         \item Bounded invertibility of closed operators and bounded operators. 
%         \item His definition of the spectrum. I spent time trying to compare this to the classical functional analysis definition in Kreyszig's text. 
%         \item Spectrum of infinite direct sums of operators. 
%         \item Spectrum of closed and densely defined operators. 
%         \item Symmetric vs. self-adjoint vs. essentially self-adjoint operators. 
%     \end{itemize}
% \end{enumerate}


\begin{enumerate}
    \item The repository is now up and running, with meeting logs and daily progress logs.

    \item Last week, I focused mainly on reviewing functional analysis. This included:
    \begin{itemize}
        \item reviewing linear, bounded, closed, compact, and inverse operators;
        \item reviewing Riesz's lemma, the Open mapping theorem, the Closed graph theorem, the Bounded Inverse theorem, Neumann series, Riesz's representation theorem, and the Riesz map. 
        \item studying the spectrum of a linear operator, including the distinction between point, continuous, and residual spectra;
        \item learning about Banach and Hilbert adjoints from Umberto's first lecture;
        \item comparing self-adjoint, unitary, and normal operators;
        \item learning about the Fredholm alternative for compact linear operators from Umberto's first lecture.
    \end{itemize}

    \item I also went through the first chapter of Nick's text. This included:
    \begin{itemize}
        \item the general motivation for computing eigenvalues;
        \item the motivation for computing pseudospectra of nonnormal operators;
        \item eigenvalues of matrices and linear operators;
        \item pseudospectra of matrices;
        \item pseudospectra of linear operators.
    \end{itemize}

    \item To further consolidate this material, I went through Umberto's first two lectures. Here, I learned about:
    \begin{itemize}
        \item Hilbert, Banach, and formal adjoints;
        \item spectra and pseudospectra of linear operators;
        \item spectral radii of bounded operators;
        \item properties of spectra for different classes of operators. For example, the spectrum of a compact operator consists of isolated points.
        \item separable Hilbert spaces and Hilbert bases;
        \item the spectral theorem for operators that are compact and self-adjoint or compact and normal.
    \end{itemize}

    \item I went through the preliminary section of Matt's text to fill in any remaining gaps. Here, I learned more about:
    \begin{itemize}
        \item separable Hilbert spaces and orthogonal projections;
        \item the Uniform Boundedness Theorem;
        \item Matt's definition of the adjoint;
        \item boundedly invertible operators;
        \item bounded invertibility of closed and bounded operators;
        \item Matt's definition of the spectrum, and how it compares with the classical functional analysis definition in Kreyszig's text;
        \item spectra of infinite direct sums of operators;
        \item spectra of closed and densely defined operators;
        \item symmetric, self-adjoint, and essentially self-adjoint operators.
    \end{itemize}

    \item Finally, I started to type up some of my handwritten notes on what I had learned. I think this is helping me connect and consolidate the material. I want to continue this habit!
\end{enumerate}


\subsection{Meeting Notes}

\begin{itemize}
    \item In the repository, I should upload the \texttt{.tex} files instead of the compiled PDFs.

    \item In my background notes, I should cite specific theorems using a format such as \verb|\cite[Theorem 1]{person}.|

    \item After finishing my written notes, I should move on to the finite element discretization of eigenproblems. The main source for this will be Boffi's text.

    \item I should also start working through examples of eigensolves in Firedrake, possibly using tutorials. Eventually, I should work through the curl--curl example to see the potential issues that arise when computing spectra.

    \item After learning about these discretizations, it may be a good time to go through the literature on multilevel solvers for eigenvalue problems. It would be useful to understand the current state of the art.

    \item Firedrake assembles PETSc vectors and matrices, and uses PETSc's KSP objects for linear solves and SNES objects for nonlinear solves. On top of PETSc is SLEPc, the ``Scalable Library for Eigenvalue Problem Computations." SLEPc adds an EPS object for eigenproblems. Thus, Firedrake's linear eigensolve creates an EPS object. The SLEPc manual and technical notes are useful resources for learning about eigensolves and their implementation.

    \item The next meeting is Monday, May 25, at 3 p.m. UK time.

    \item I have a meeting with Matt Colbrook sometime on June 11.

    \item The presentation is in the first week of June. I will start thinking about what this should look like soon.
\end{itemize}


\subsection{This Week's Plan}

\begin{enumerate}
    \item Today (May 19) and tomorrow, I will work on finishing up my background notes. I hope to get this done!
    \item I will then move on to learning about finite element discretizations of eigenproblems. This will involve some general googling, Boffi's text, and maybe the SLEPSc manual. 
    \item Work through some examples in Firedrake. First, I will work through some tutorials. I hope to work up to the curl-curl example. 
    \item Transition to a literature review for Multilevel eigensolves.
\end{enumerate}
% \begin{itemize}
%     \item In the repo, I should upload .tex files instead of the compiled pdf's.
%     \item In my background notes, cite the specific theorems using "[Theorem 1]\cite{person}."
%     \item After finishing up my written notes, I should move on to the finite element discretization of eigenproblems. The main source for this will be Boffi's text. 
%     \item I should also start doing some examples of eigensolves in Firedrake. I could do some tutorials. Also, eventually, I should do this curl-curl example to actually see the potential issues in computing spectra.
%     \item Then, after learning about these discretizations, it may be a good time to go through the literature on multilevel solvers for eigenvalue problems. It would be good to understand the current state of the art.
%     \item Firedrake assembles PETSc vectors and matrices, and uses the ksp (linear solve) and snes (nonlinear solve) PETSc commands to solve the discretized problems. On top of PETSc is SLEPc (a scalable library for eigenvalue problems). SLEPSc adds an eps object used for eigenproblems. Thus, Firedrake's linear eigensolve creates an eps object. The SlEPSc manual and technotes are a great resource for learning about eigensolves and their implementation!
%     \item Next meeting is Monday, May 25, at 3 pm UK time. 
%     \item Meeting with Matt Colbrook sometime on June $11^{th}.$
%     \item Presentation first week of June. I will start thinking about what this looks like soon. 
% \end{itemize}


\section{Daily Logs}

\textbf{Tuesday} 

\begin{itemize}
    \item Spent the morning prepping for the meeting with Patrick, having the meeting with Patrick, and creating a new meeting log.
    \item I then returned to my typed-up background notes. I made the citation changes suggested by Patrick. I then proceeded to start adding in the proofs I had left out. I will continue with this tomorrow and the next day. 
\end{itemize}

\textbf{Wednesday} 

\begin{itemize}

    \item Today was a travel day. Unformtunately, between commuting, flights, etc, not much got done. For a couple of hours on the plane I worked on typing up some proofs.
\end{itemize}



\textbf{Thursday} 

\begin{itemize}
    \item Today was my family's event. This meant that I only got about ffive-ish hours of work done. I proceeded to finish up some typed up proofs. 
    \item I will move on to the finite element background tomorrow!
\end{itemize}

\textbf{Friday}

\begin{itemize}
    \item I spent the  first two hours cleaning up some final proofs in my written notes.
    \item Then I moved on to Boffi's text. I did a little FEM and Firedrake review to set the scene. Then I proceeded to write notes on the the simplest example (1D Laplace) and coded it up in Firedrake. 
\end{itemize}
    
\textbf{Sunday}
\begin{itemize}
    \item Working on 2D Laplace examples today.
    \item Took a while to figure out how to create unstructured meshes. used Chat to help. Not sure if I should use netgen or gmsh for the unstructured square meshes. Not sure how to pick a desired number of DOF in netgen. How did Boffi construct their meshes? 
    \item Worked through some of the Laplace eigenvalue examples in 2D. In particular, I worked on unstructured meshes, structured meshes, meshes with reverse diagonals, meshes with criss-crosses, and unstructured L-shaped meshes.
    \item I then did a quick review on conforming vs. nonconforming elements and recreated the Boffi example with Crouzeix-Raviart elements. As expected, the eigenvalues are now approximated from below. 
    \item Tomorrow (and the next day), I plan to work through the examples for the mixed formulation of the Laplace eigenvalue problem. Then, finally, I will move on to the Maxwell problem. 
    \item What sections of Boffi should I work through and study? All of them? 
    \item At some point, I want to spend some time looking through SLEPc to try and understand the solvers some more. 
\end{itemize}

\end{document}
